\subsection{Denominators of the projections on types}
\hypertarget{source-graded-projectors}{}

The preceding Gram calculation gives a rational decomposition. Its use in integral interpolation also requires the denominators of the projections to the graded types. In the
first, let \(E_{ij}\) be the matrix units acting on a
finite-dimensional rational \(\mathrm{GL}_r\)-module in
characteristic zero. Suppose an integral lattice respects the
weight decomposition and is stable under
\(E_{ij}^{a}/a!\) for all \(a\) and \(i\ne j\). We use complete reducibility, highest weights and the Pieri
and branching rules for split general linear groups in
characteristic zero. Their Schur-module constructions
commute with extension of the ground field to
\(\mathbb Q_p\). \cite{fulton-harris}
The polynomial
lattices in the letters have this property by taking the
coefficients of the actions of elementary unipotent matrices.
The lattice of homomorphisms between two such modules has it
as well, by expanding the left and inverse right unipotent
actions.

\begin{lemma}[Denominators for invariant and isotypic projections]
\label{lem:integral-type-projection}
If the weights in the module are in \([-l,l]\) in each coordinate and
\(l+r-1\ge1\), write \(P\) for projection onto the invariants and \(M_0\) for the zero-weight lattice. The denominator bound is
\[
 p^{\binom r2\lfloor\log_p(l+r-1)\rfloor}\,P(M_0)\subseteq M_0.
\]
 The same multiplier
works for an isotypic projector when \(l\) bounds the coordinate weights
of the corresponding Hom representation.
\end{lemma}

\begin{proof}
\proofstep{Projecting onto invariants by the integral trace.}
On the zero weight space one can project onto the invariant
vectors with denominator at most
\[
                  p^{\binom r2\lfloor\log_p(l+r-1)\rfloor},
\]
when all coordinate weights in the module lie in \([-l,l]\).
Take \(l+r-1\ge1\). To see this, inductively suppose that
the input \(v\) is invariant in a set \(I\) of \(s\) coordinates.
Adjoin coordinate \(z\). Put \(\Omega_0=1\),
\(\Omega_{q+1}=\sum_{i\in I}E_{zi}\Omega_q E_{iz}\),
and apply
\[
   P_{I,z}=\sum_{q=0}^l c_q\Omega_q,\qquad
   c_q=\frac{(-1)^q s!}{(q+s)!q!}.
\]
Each \(\Omega_q\) commutes with \(\mathfrak{gl}_I\) by contracting
the two indices in the recursion, and preserves weights.
For \(j\in I\), commuting the matrix units gives
\[
 E_{jz}\Omega_q v=\Omega_q E_{jz}v+
                       q(q+s)\Omega_{q-1}E_{jz}v
                 \qquad(q\ge1).
\]
For example expand \(\Omega_q\) as the sum of all products of
\(q\) factors \(E_{zi}\) followed in reverse order by the
corresponding \(E_{iz}\).

In moving \(E_{jz}\) through the
first \(q\) factors each commutator is
\(E_{ji}-\delta_{ij}E_{zz}\). Moving these through the remaining
factors and using \(E_{ii'}v=E_{zz}v=0\) for \(i,i'\in I\)
gives the stated recurrence. Substituting the coefficient at the next index gives, for \(0\le q<l\),
\[
 (q+1)(q+s+1)c_{q+1}
 =(q+1)(q+s+1)\frac{(-1)^{q+1}s!}{(q+s+1)!(q+1)!}=-c_q .
\]
Thus the two contributions at each interior index cancel.
Its endpoint \(\Omega_l E_{jz}v\) vanishes, because the inner
\(l+1\) raising factors move the \(z\)-weight below \(-l\).
Thus every \(E_{jz}\) kills the result. On its zero weight
the opposite \(E_{zj}\) also kills it, by the finite-dimensional
\(\mathfrak{sl}_2\) string for \((j,z)\). Vectors already
invariant are fixed. All terms beyond the identity are in the
span of images of off-diagonal generators.

For a multi-index \(a=(a_i)_{i\in I}\) with \(\sum a_i=q\), expand
\(c_q\Omega_q\) using commutation within each of the two groups
of root operators. The coefficient of the product of their
\emph{divided} powers has absolute rational value
\(s!\prod_i a_i!/(s+q)!\).
Number the \(s\) indices in order. Cancelling successive factorials gives
\[
 \frac{(s+q)!}{s!\prod_i a_i!}=\prod_{j=1}^s\binom{s+a_1+\cdots+a_j}{a_j}.
\]
Each has valuation at most \(\lfloor\log_p(l+s)\rfloor\)
by the carry bound. Summing the denominator exponents for
\(s=0,\ldots,r-1\) proves the claim. The construction commutes
with intertwiners, since it uses only weight projection and the
matrix units. On an endomorphism space it preserves the trace.
Projection onto weight zero retains all diagonal entries, and
every subsequent correction is a sum of commutators.

\proofstep{Denominators on a highest-weight summand.}
In particular take an irreducible summand \(U\) and the module
\(\operatorname{Hom}(U,V)\), with action
\(E\phi=E_V\phi-\phi E_U\). Use the coordinate weight
bound \(l\) on differences of weights. Denote the invariant
projection just proved, preceded by weight-zero projection,
by \(P\). With an integral basis \(u_j\) of a divided-power
stable weight lattice in \(U\) the operator
\[
     Q_U(v)=(\dim U)\sum_j P(u_j^*\otimes v)(u_j)
\]
is the isotypic projector. Such a lattice is obtained by
intersecting a copy of the summand defined over \(\mathbb Q_p\) with the ambient integral
lattice. Its preservation under the divided powers follows by
intersection. On an inequivalent irreducible the invariant Hom
space is zero.

For an embedding \(i:U\to V\), naturality and
Schur's lemma give
\[
 P(u_j^*\otimes i(u))=i\circ(u_j^*(u)/\dim U)\operatorname{id}_U,
\]
since the invariant projection preserves trace on endomorphisms.
This gives \(Q_U(i(u))=i(u)\). Hence the same denominator bound
works for a sum of any chosen types. For letters of fixed total
degree \(S\), the span of coordinate weights in the full module
is at most \(S\): a vector of degree \(d\) contributes in \([0,d]\), and a covector of degree \(d\) in \([-d,0]\). The width of the total interval is therefore bounded by
\[
 \sum_{\rm vectors} d+\sum_{\rm covectors} d=S.
\]
Thus one may use \(l=S\) also for the relevant Hom spaces.
In the Fock metric \(E_{ij}\) is adjoint to \(E_{ji}\), including
the contragredient action on covectors, so distinct types are
orthogonal.

For the second estimate, fix an order \(e\) and let \(M_e\) be
the lattice of initial terms \(I_e f\). Here the integral \(f\)
are in a fixed full weight, highest for the coordinates other
than \(\rho\), and satisfy the prefix bounds on first coordinates
of the generated modules. Put \(H_e=\mathbb Q_pM_e\) and
\[
 L_e=H_e\cap\mathbb Z_p[\mathbf x].
\]
Suppose a self-adjoint
projection onto \(H_e\) has denominator at most \(p^d\) on
the monomials of the required multidegree and that
\(p^c L_e\subseteq M_e\). The graded determinant allowance
is then \(v_*=d+2c\). Indeed a primitive \(f\in M_e\) has a
coefficient of valuation \(\le c\). Otherwise
\[
 f\in p^{c+1}L_e\quad\Longrightarrow\quad f/p\in p^cL_e\subseteq M_e,
\]
contradicting primitivity.

If its exponent is \(\alpha\),
the normalized Fock pairing with \(\mathbf x^\alpha\) has valuation
at most \(c\), because the reciprocal of the monomial norm is a
product of integral multinomial coefficients. Multiply the
projected monomial by \(p^{d+c}\) to obtain an integral partner
in \(M_e\) with pairing of valuation \(\le d+2c\).
Using the row corresponding to a deepest invariant factor in
a Smith decomposition, this primitive-vector test bounds every
invariant factor by \(d+2c\), and thus the determinant valuation
by \((d+2c)\operatorname{rank}M_e\). The zero-rank convention has
determinant one.

\end{proof}

\subsection{Lifting integral initial terms}

Clearing a projector controls its image. We also need to lift prescribed leading terms into that image while keeping the prefix bounds.
Assume \(r\ge2\), take \(\rho\) last, and let \(K\) be the number
of letters preceding the final vector \(w\). A cutoff on a prefix
means that every irreducible type of that prefix occurring in
the vector has first highest-weight coordinate at most the specified
integer. If \(C\) is the cutoff and \(G\) the prefix group, the test on a nonzero polynomial \(g\) is
\[
 \max\{\mu_1:(xg)_\mu\ne0\text{ for some }x\in G\}\le C,
\]
where the subscript denotes a weight component. An infinite cutoff imposes
no condition. If \(f=w_\rho^e h\), a prefix containing \(w\)
satisfies its cutoff precisely when \(h\) satisfies the cutoff
lowered by \(e\). Indeed for a generic group translate the top
first coordinate in the translated linear form \(w_\rho\) is one.

Let \(m_1(f)\) be the largest first coordinate weight among
prefix translates of \(f\), and similarly for \(h\). On a common
Zariski-open set the highest first weight of the translated
\(w_\rho\) is one, and the component of \(h\) realizing \(m_1(h)\)
is nonzero. Their top components multiply to a nonzero polynomial,
so
\[
                  m_1(w_\rho^e h)=e+m_1(h).
\] For shorter prefixes the
factor is fixed. The group on the coordinates other than \(\rho\)
also fixes it. We may therefore divide out this power, subtract
\(e\) from the last degree and the \(\rho\)-weight, and take
restriction \(R\) at \(w_\rho=0\). Multiplication back by
\(w_\rho^e\) recovers the graded term. It scales the Fock
form on the restricted monomials by a common factor.

\begin{lemma}[Saturation of the restricted coefficient map]
\label{lem:graded-coefficient-lift}
For \(r\ge2\), with the prefix-first-weight cutoffs and
restriction at \(w_\rho=0\), the image of restriction
on the integral polynomial lattice contains the integral restricted
lattice after the multiplier in \eqref{eq:res-sat} before the highest
weight equations. After the weight and type projectors let \(d_{\rm gr}\) be the normalized valuation per rank of the graded determinant. With \(K\) the number of earlier letters,
\[
 d_{\rm gr}\le(K+3)r^2(1+\log_p(8a+r))\qquad(S\le2a).
\]

\end{lemma}

\begin{proof}
\proofstep{The lowered last degree.}
Keep the shortened prefix cutoffs fixed and vary the reduced last
degree \(t\). Write \(Y_t\) for the space satisfying these cutoffs
before restriction, with no highest-vector or weight requirements.
In the restricted space define \(X_t\) by the same earlier prefix
conditions. If there is a finite full-prefix cutoff, let \(C\) be
the least one. Impose in addition the bound \(C\) on the first
\(K\) letters under \(\mathrm{GL}_r\) and on the entire restricted
word under \(\mathrm{GL}_{r-1}\).

We have \(RY_t\subset X_t\).
Indeed, taking a prefix type and tensoring it by the last symmetric
power adds a horizontal strip, by the Pieri rule (twist by a
determinant to reduce integral highest weights to partitions).
If \(\alpha\) is the earlier type and \(\lambda\) a resulting type, the strip condition gives
\[
 \lambda_1\ge\alpha_1\ge\lambda_2.
\]
In particular a cutoff on \(\lambda_1\) also bounds \(\alpha_1\). Restriction commutes with the earlier
group, and with \(\mathrm{GL}_{r-1}\) on the whole word.
Interlacing bounds the first coordinate of the latter subgroup
types by that of the full type.

Let \(S\) be the total degree before the shortening, and set
\[
 A_0=\binom r2\lfloor\log_p(S+r-1)\rfloor,\qquad
 A_1=\binom r2\lfloor\log_p(4S+r-1)\rfloor .
\]
We claim that at the required last degree \(t\)
\begin{equation}\label{eq:res-sat}
  p^{2A_1}(X_t\cap\mathbb Z_p[\mathbf x])
                  \subset R(Y_t\cap\mathbb Z_p[\mathbf x]).
\end{equation}
With no full cutoff, or with \(C\) at least the total reduced
degree at this \(t\), the polynomial independent of \(w_\rho\)
lifts itself. For the other case use an auxiliary Lie algebra action.
Denote the surviving coordinates by \(I\), put \(E_w=\sum w_i
\partial_{w_i}\), and write \(B_{ab}\) for the matrix units on
the first \(K\) letters. Use an extra index \(*\) preceding all
of \(I\). The following operators act on \(\bigoplus_t X_t\):
\begin{align*}
 G_{ab}&=B_{ab}+w_a\partial_{w_b},&
 G_{**}&=C-E_w,\\
 G_{*a}&=\partial_{w_a},&
 G_{a*}&=w_a(C-E_w)-\sum_{b\in I} w_b B_{ab}.
\end{align*}
The relations \([\partial_i,w_j]=\delta_{ij}\) and the matrix-unit brackets for \(B\) give
\[
 \begin{aligned} {}[G_{ab},G_{c*}]&=\delta_{bc}G_{a*},\\ {}[G_{*b},G_{a*}]&=\delta_{ab}G_{**}-G_{ab},\\ {}[G_{a*},G_{b*}]&=0.\end{aligned}
\]
Together with the degrees these give the
\(\mathfrak{gl}_r\) brackets.

\proofstep{Preservation of the endpoint cutoff.}
We check the cutoff at the endpoint, where raising could otherwise
leave the space. The same formulas apply on all coordinates before
restriction, with the sum over all \(b\); restriction intertwines
them for \(a,b\in I\), since every omitted summand has a factor
\(w_\rho\). At an earlier prefix each \(B_{ab}\) splits into an
operator of the prefix group and one on disjoint letters.
Both preserve its type condition. At the full prefix the raises
transform as a vector under the coordinate group and the derivatives
as a covector. A derivative never increases first coordinate
weight. A raise with index different from the first has first
weight zero.

For the first index, denoted by 1, the identity
\[
                   G_{1*}=C w_1-\sum_b w_bG_{1b}
\]
shows that it kills the component of first weight \(C\).
On this component \(G_{11}=C\) and each \(G_{1b}\) with
\(b\ne1\) vanishes because it raises the first weight beyond
the cutoff. On lower weights the result stays below or at \(C\).
Applying this argument after every group translate proves the
stable type bound for all the raises, upstairs and downstairs.

The divided powers also preserve integral coefficients. This
is immediate for differentiation. For the raise with index \(a\),
let \(w\) be the column of variables, and let \(\mathcal B\)
be the representation on the earlier letters. On a polynomial
\(h\) valued in that representation the generating series is
\[
 (1+u w_a)^C\,
 \mathcal B\!\left(I-\frac{u e_a w^{\mathsf T}}{1+u w_a}\right)
                 h\!\left(\frac{w}{1+u w_a}\right).
\]
Let \(T(u)\) denote this substitution. The rank-one matrix product
gives \(T(u+v)=T(u)T(v)\), and \(T(0)=1,\ T'(0)=G_{a*}\).

Differentiate in \(v\) at zero and compare coefficients in \(u\):
\[
 T'(u)=T(u)G_{a*},\qquad
 (j+1)[u^{j+1}]T=([u^j]T)G_{a*},\qquad
 [u^j]T=G_{a*}^j/j! .
\] The displayed matrix and its inverse
\(I+u e_aw^{\mathsf T}\) have integral series entries, so this
also works for covectors and negative \(C\). For the operators on
\(I\) the ordinary elementary matrices give integral divided
powers. Thus the integral lattice in the direct sum of the \(X_t\),
and the sum of the images of the integral \(Y_t\), are stable
under divided powers and their integral weight decompositions.

\proofstep{Lifting the initial forms over the coefficient ring.}
For completeness this is a finite-dimensional calculation. Let
\(A\) be the total earlier degree. In a nonzero \(X_t\) a generic
\(\mathrm{GL}_{r-1}\) translate has a nonzero coefficient of the
\(t\)-th power of the first \(w\)-coordinate. Indeed a nonzero
homogeneous vector-valued polynomial is nonzero at some point,
and the simultaneous change of the earlier letters is invertible.
Its earlier letters have first weight at least \(-A\). The cutoff at \(C\) therefore gives
\[
 t-A\le C,\qquad t\le C+A,\qquad C-t\ge -A.
\]
All coordinate weights lie between \(-A\) and \(C<S\). In particular the coordinate weight differences
in the Hom construction are bounded by \(4S\). Highest vectors
for this auxiliary ordering are killed by every \(G_{*a}\).
In characteristic zero they are independent of \(w\). At degree
zero the early cutoff \(C\) is already the full cutoff, so all
integral highest vectors lift identically.

To see the index conclusion, first apply the type projectors to
an integral element of the direct sum, multiplying by \(p^{A_1}\)
to make each component integral. For a component of type \(U\)
use in \(U\) the lattice generated by divided powers from an
integral highest vector. This is a lattice: it spans by highest
weight theory and lies in a fixed stable lattice after scaling
the highest vector. Let \(L_{\rm im}\) be the integral image of the lift
and
\[
 T_j=p^{A_1}P(u_j^*\otimes v).
\]
These integral
intertwiners send the highest generator into \(L_{\rm im}\).
Every \(u_j\) is generated from that highest vector by
divided powers, which also preserve \(L_{\rm im}\).
Hence \(T_j(u_j)\in L_{\rm im}\) for every \(j\), and
the isotypic projector formula reads
\[
       p^{A_1}v=(\dim U)\sum_j T_j(u_j)\in L_{\rm im}
\]
on the component of type \(U\). Sum over the components and retain the required
homogeneous last degree of a lift to obtain \eqref{eq:res-sat}.

\proofstep{Imposing the highest-weight equations.}
It remains to impose the weight and highest-vector conditions.
Full weight projection is integral. Within that weight the
highest vectors for \(I\) are precisely the copies of the
isotypic summand with highest weight equal to the specified
\(I\)-weight, intersected with the weight space. This projector
costs at most \(A_0\), commutes with restriction and preserves
every upstairs prefix condition. Indeed those conditions are
stable under the total action on \(I\) and the full torus.
Thus the saturation exponent in the primitive-vector test is
at most
\[
 c=2A_1+A_0.
\]
For its self-adjoint projector first
retain monomials of the required \(w_\rho\)-order and weight.
On the restricted polynomials impose the early prefix type
projectors including the cutoff at \(K\). They commute: a
projector on an earlier prefix commutes with its group action
and with all operators on later letters, hence with the action
on the longer prefix.

The prefix projectors also commute with the total \(I\)-type projector. That projector selects the required
highest type, giving zero if it violates the full cutoff.
Repeated bounds on the same prefix take their minimum. The
zero-letter and one-letter prefix tests are scalar on a fixed
multidegree. For \(K\ge1\) at most \(K+1\) projectors with
denominator \(A_0\) suffice. These self-adjoint projections have
the required rational image by \eqref{eq:res-sat}; putting back
\(w_\rho^e\) preserves self-adjointness. We obtain
\[
             v_*\le (K+3)A_0+4A_1
             \le (K+3)r^2(1+\log_p(8a+r))
                        \qquad(S\le2a,\ K\ge1).
\]
The last inequality uses \(K+7\le2(K+3)\). For \(K=0\) the
original single symmetric power is irreducible, so the cutoff
is scalar and every allowed restricted polynomial lifts.
The weight and highest tests alone suffice and obey the same
final bound. For \(e\) greater than the last degree the graded
piece is zero.
\end{proof}
